\chapter{Giữ Nhịp Giữa Tiết}
\chapterintro{Đừng đợi lớp rơi hẳn mới cứu}{Giữa tiết học, chỉ cần đổi nhịp đúng lúc thì lớp sẽ không bị trôi xa.}

\section{Khi năm bạn bắt đầu nhìn ra cửa}
\begin{tinhhuong}{Tình huống}
Thầy cô đang giảng bình thường nhưng vài em đầu tiên đã phân tán. Nếu bỏ qua, phần còn lại sẽ đi theo rất nhanh.
\end{tinhhuong}
\begin{goiydungngay}
Ngắt mạch bằng câu: \textit{``Nếu phải giữ lại đúng một ý trong 10 giây qua, em giữ ý nào?''}
\begin{itemize}
\item Chỉ gọi một em bất kỳ.
\item Không chê câu trả lời ngắn.
\item Sau khi nghe, tóm lại bài bằng đúng một câu.
\end{itemize}
\end{goiydungngay}
\begin{cauchot}
Muốn giữ lớp, hãy buộc suy nghĩ quay lại trước khi mắt quay đi quá xa.
\end{cauchot}
\ngancach

\section{Cả lớp đứng lên chọn góc}
\begin{tinhhuong}{Tình huống}
Lớp bắt đầu ì, ai cũng ngồi lâu, không khí xuống rõ rệt. Cần một thay đổi cơ thể thật ngắn.
\end{tinhhuong}
\begin{goiydungngay}
Đưa ra hai lựa chọn đơn giản: \textit{``Đồng ý đứng góc trái, còn phân vân đứng góc phải.''}
\begin{itemize}
\item Dùng với câu hỏi nhẹ, không quá nhạy cảm.
\item Cho học sinh đổi góc nếu nghe bạn nói hợp lý.
\item Sau 90 giây, mời lớp ngồi xuống và chốt bài.
\end{itemize}
\end{goiydungngay}
\begin{cauchot}
Đổi vị trí cơ thể rất ngắn có thể kéo lại vị trí của sự chú ý.
\end{cauchot}
\ngancach

\section{Hỏi ngược thay vì giảng tiếp}
\begin{tinhhuong}{Tình huống}
Thầy cô có cảm giác mình đang nói nhiều hơn mức lớp cần nghe. Không khí chưa hẳn xấu, nhưng đang dần phẳng.
\end{tinhhuong}
\begin{goiydungngay}
Dừng lại và hỏi: \textit{``Nếu em phải hỏi ngược lại thầy cô một câu, em sẽ hỏi gì?''}
\begin{itemize}
\item Câu hỏi ngược thường làm lớp tỉnh hơn.
\item Có thể cho học sinh hỏi bằng giấy nếu ngại nói.
\item Dùng 1 đến 2 câu hỏi thôi, rồi quay lại bài.
\end{itemize}
\end{goiydungngay}
\begin{cauchot}
Có lúc cứu lớp không phải bằng thêm lời giảng, mà bằng một câu hỏi trả quyền chủ động lại cho học sinh.
\end{cauchot}
\ngancach

\section{Nhắc lại bằng tay trước khi nhắc lại bằng miệng}
\begin{tinhhuong}{Tình huống}
Giữa tiết, lớp chưa buồn ngủ hẳn nhưng đã có dấu hiệu nghe mà không nhập.
\end{tinhhuong}
\begin{goiydungngay}
Mời lớp dùng tay làm nhanh một tín hiệu: giơ một ngón cho ý rõ nhất, hai ngón cho chỗ còn lăn tăn.
\begin{itemize}
\item Không mất quá 10 giây.
\item Thầy cô nhìn là biết nên chậm lại hay đi tiếp.
\item Sau đó chỉ gọi đúng 1 em nói thêm.
\end{itemize}
\end{goiydungngay}
\begin{cauchot}
Một tín hiệu nhỏ bằng tay nhiều khi cứu nhịp nhanh hơn một đoạn hỏi đáp dài.
\end{cauchot}
\ngancach

\section{Đổi sang một câu dự đoán}
\begin{tinhhuong}{Tình huống}
Phần giảng đang đều nhưng bắt đầu phẳng, lớp nghe mà mắt không còn sáng.
\end{tinhhuong}
\begin{goiydungngay}
Dừng lại và hỏi: \textit{``Theo em, đoạn tiếp theo dễ làm lớp nhầm ở chỗ nào?''}
\begin{itemize}
\item Học sinh thích dự đoán hơn là chỉ nghe tiếp.
\item Câu hỏi này kéo sự chú ý về phần sắp tới.
\item Nghe 2 ý rồi mới giảng tiếp.
\end{itemize}
\end{goiydungngay}
\begin{cauchot}
Dự đoán làm lớp đi trước một bước, và chính bước đó giữ các em ở lại với bài.
\end{cauchot}
